\begin{textsample}{POS Dim 3 – al – Score 13.00 – al/al\_inf\_16.txt}  \label{ex:f3_pos_176}
Desenvolvimento e aprendizagem e os grupos etários Quando se fala em educação é muito comum utilizar a palavra aprender, sendo esta geralmente associada a conteúdo, conhecimento, mais voltado a práticas escolarizantes.
No entanto, é importante destacar que na educação \textbf{infantil} a palavra aprendizagem sempre vem acompanhada da palavra desenvolvimento, sendo este voltado para todas as dimensões humanas.
Atuar na educação \textbf{infantil} exige que o professor se aproprie de conhecimentos específicos sobre o desenvolvimento \textbf{infantil} em sua totalidade e compreenda que o aprendizado se refere à aquisição de diversas habilidades ligadas a este desenvolvimento.
Pensar um currículo que alcance o objetivo da educação \textbf{infantil} na prática, desenvolvimento integral da \textbf{criança}, se faz necessário conhecer as especificidades de cada etapa do desenvolvimento destes sujeitos.
Para definir referenciais sobre desenvolvimento \textbf{infantil} teremos como referencial a teoria do de Henry Wallon, sendo este um aporte teórico que contempla todas as fases e dimensões do desenvolvimento da pessoa humana.
Deste modo, exibiremos algumas características por grupos etários conforme organizados pela Base : \textbf{bebês} ( 0 a 1 ano 6 \textbf{meses} ), \textbf{crianças} bem \textbf{pequenas} ( 1 ano e 7 \textbf{meses} a 3 anos e 11 \textbf{meses} ) e \textbf{crianças} \textbf{pequenas} ( 4 anos a 5 anos e 11 \textbf{meses} ). - \textbf{Bebês} ( zero a 1 ano e 6 \textbf{meses} ) Nesta primeira fase do desenvolvimento compreende dois momentos, que vai do nascimento ate 1 ano de idade, denominado por Wallon como estágio impulsivo-emocional.
Inicialmente a \textbf{criança} vive em total dependência do meio externo.
A \textbf{criança} vai se descobrindo enquanto sujeito que interage e estabelece relações emocionais, a partir de sensações e emoções, por isso, a necessidade do toque.
Nessa fase, geralmente as ações não são intencionais, a aprendizagem ocorre “ acidentalmente ”, por reflexos. “ Paulatinamente essa interação criança-meio externo construirá um campo de comunicação recíproca ”. ( MAHONEY et al, 2007, p.19 ).
Para tanto, dispor de um espaço com objetos coloridos e diversificados é fundamental para estimular a interação entre os \textbf{bebês}, com os objetos, com os \textbf{adultos} e com o ambiente.
Sendo assim, atividades com músicas, comandos de voz, conversa com os \textbf{bebês} é essencial para o essa fase.
Promover situações nas quais a \textbf{crianças} experimente sensações que estimulem movimentos favorecendo a comunicação, a maturação do tônus muscular, a afetividade, consequentemente os processos psíquicos.
Com algumas habilidades preliminares desenvolvidas, na segunda metade do primeiro ano de vida, a \textbf{criança} começa a realizar atividades repetidas que preludiam o estágio posterior : o sensório-motor.
Estas atividades são denominadas circulares que : > [... ] são movimentos inicialmente casuais, mas que serão repetidos intencionalmente pela \textbf{criança}, levando -a a investigar a conexão entre seus movimentos e seus efeitos e a variação dos efeitos diante das variações de movimentos, ajustando cada vez mais seus \textbf{gestos} aos resultados obtidos e tomando -os, assim, mais precisos e úteis.
Tais atividades levam a \textbf{criança} a conhecer cada vez mais a si própria e aos objetos, e acontecem em função do prazer pela repetição, da perseverança indispensável ao processo de aprendizagem, e da motivação investigadora presente nas \textbf{crianças} para descobrir as conexões entre seus atos e os respectivos efeitos ” ( MAHONEY et al, 2007, p.27 ).
O professor precisa estar atento à comunicação estabelecida entre eles, como estão interagindo uns com os outros, pois em nenhum outro período da vida o ser humano faz tantas conquistas motoras, mentais e sociais quanto nos primeiros \textbf{meses} vida dos sujeitos.
As atividades didático-pedagógicas devem ser planejadas conforme as especificidades do desenvolvimento de cada \textbf{bebês} com sequências lógicas, pois, cada habilidade nova devem ser aperfeiçoadas conforme uma já aprendida.
A ação didática deve promover diferentes situações de exploração do corpo para estimular as musculaturas dos olhos, em seguida sustentar o pescoço, o tórax, até ficar em \textbf{pé}.
Nessa fase, o olhar do \textbf{bebê} é atraído por objetos movimento e de cores contrastantes, como preto e branco.
As \textbf{paredes} das escolas devem ser atrativas com cores fortes, a organização dos móveis, de forma que promova a interação entre eles e com as pessoas.
Os \textbf{brinquedos} devem ser acessíveis para que o \textbf{bebê} alcance os objetos que deseja.
Também chuta, se \textbf{balança}, se debate e \textbf{bate}, esfrega, arranha, se inclina de modo rítmico e repetitivo.
Com isso, manda estímulos para o cérebro, que se organizam em informações para as próximas habilidades : engatinhar, ficar em \textbf{pé} e andar.
A música é um veículo que estimula a comunicação e expressão, além de favorecer o prazer pelos \textbf{sons}.
Além de atividades direcionadas, podem acontecer “ cantinhos de músicas ” com instrumentos confeccionados ou adquiridos, onde as \textbf{crianças} possam fazer experiências com \textbf{sons} de maneira independentes.

% matched lemmas: adulto, balançar, bater, bebê, brinquedo, criança, gesto, infantil, mês, parede, pequeno, pé, som
\end{textsample}
